\section*{Chapitre \ref{ch:g11:seismic-structure} --- Les ondes sismiques révèlent l'intérieur}

\begin{solution}{exo:g11:seismic-structure:1}
Une \omterm{def:g11:seismic-structure:body-waves}{onde P} comprime et étire la roche, qui se déplace dans la direction de
propagation~; une \omterm{def:g11:seismic-structure:body-waves}{onde S} la cisaille, et la roche se déplace
perpendiculairement à cette direction. Les \omterm{def:g11:seismic-structure:body-waves}{ondes P} traversent les solides et
les liquides~; les \omterm{def:g11:seismic-structure:body-waves}{ondes S} ne traversent que les solides. L'\omterm{def:g11:seismic-structure:body-waves}{onde P} est
enregistrée la première, parce que c'est la plus rapide des deux.
\end{solution}

\begin{solution}{exo:g11:seismic-structure:2}
$40 \times 111.2 = \qty{4448}{km}$, environ \qty{4450}{km}. $8900/111.2
\approx \ang{80}$.
\end{solution}

\begin{solution}{exo:g11:seismic-structure:3}
\qty{456}{s} font 7~minutes 36~secondes~; \qty{822}{s} font 13~minutes
42~secondes. Le délai S--P vaut $822 - 456 = \qty{366}{s}$, soit 6~minutes
6~secondes.
\end{solution}

\begin{solution}{exo:g11:seismic-structure:4}
$\sin i_c = 6.0/8.0 = 0.75$, donc $i_c \approx \ang{48.6}$.
\end{solution}

\begin{solution}{exo:g11:seismic-structure:5}
Le noyau externe a une épaisseur de $5150 - 2891 = \qty{2259}{km}$. Le rayon
du noyau vaut $6371 - 2891 = \qty{3480}{km}$, et $3480/6371 \approx 0.55$~:
un peu plus de la moitié du rayon de la Terre.
\end{solution}

\begin{solution}{exo:g11:seismic-structure:6}
6~minutes 6~secondes font \qty{366}{s}~: le séisme est à \ang{40}, soit
$40 \times 111.2 \approx \qty{4450}{km}$. Un seul délai ne donne que la
distance~: l'épicentre peut se trouver n'importe où sur un cercle de ce
rayon autour de la station. Il faut deux stations de plus pour trouver le
seul point situé à la bonne distance des trois.
\end{solution}

\begin{solution}{exo:g11:seismic-structure:7}
$i_c \approx \ang{48.6}$ et $\cos i_c \approx 0.661$. Onde directe~: $200/6.0
\approx \qty{33.3}{s}$. \omterm{ex:g11:seismic-structure:sp}{Onde conique}~: $200/8.0 + 2 \times 35 \times 0.661/6.0
\approx 25.0 + 7.7 = \qty{32.7}{s}$. L'\omterm{ex:g11:seismic-structure:sp}{onde conique} arrive la première,
environ \qty{0.6}{s} plus tôt.
\end{solution}

\begin{solution}{exo:g11:seismic-structure:8}
$\tan i_c \approx 1.134$~: la distance critique vaut $2 \times 35 \times 1.134
\approx \qty{79}{km}$. La \omterm{prop:g11:seismic-structure:crossover}{distance de croisement} vaut $2 \times 35 \times
\sqrt{14.0/2.0} = 70\sqrt 7 \approx \qty{185}{km}$. La station de
l'\cref{exo:g11:seismic-structure:7} se trouve à \qty{200}{km}, au-delà du
croisement~: l'\omterm{ex:g11:seismic-structure:sp}{onde conique} arrive la première, comme on l'a trouvé.
\end{solution}

\begin{solution}{exo:g11:seismic-structure:9}
$\sin i_c = 6.5/8.0 = 0.8125$, $i_c \approx \ang{54.3}$ et $\cos i_c \approx
0.583$. \omterm{prop:g11:seismic-structure:head-wave}{Temps d'interception}~: $2 \times 6 \times 0.583/6.5 \approx \qty{1.1}{s}$. \omterm{prop:g11:seismic-structure:crossover}{Distance de croisement}~: $12\sqrt{14.5/1.5} \approx \qty{37}{km}$. Sous un continent, le croisement se situe vers
\qty{165}{km}~: des stations espacées de quelques dizaines de kilomètres
saisissent les deux droites. Sous l'océan, l'\omterm{ex:g11:seismic-structure:sp}{onde conique} arrive déjà la
première au-delà d'environ \qty{37}{km}, et la droite de l'onde directe ne
s'étend que sur quelques dizaines de kilomètres~: il faut des stations
espacées de quelques kilomètres pour la tracer.
\end{solution}

\begin{solution}{exo:g11:seismic-structure:10}
$\sin r = (8.06/13.7)\sin\ang{30} \approx 0.294$, donc $r \approx \ang{17.1}$~:
le rai est dévié vers la perpendiculaire, c'est-à-dire vers le centre de la
Terre. Il traverse donc le noyau selon un long trajet proche de son milieu,
et ressort loin, de l'autre côté~: les rais P qui entrent dans le noyau
ressortent bien plus loin du séisme que les rais qui le frôlent.
\end{solution}

\begin{solution}{exo:g11:seismic-structure:11}
$12\,742/13.7 \approx \qty{930}{s}$, environ 15.5~minutes, alors que l'onde
met \qty{1210}{s}, environ 20~minutes~: elle a quelque 280~secondes de
retard. Même à la vitesse la plus élevée du manteau, l'onde arriverait bien
plus tôt~: le milieu de la Terre doit donc la ralentir. Oldham en conclut
que la Terre possède un noyau central dans lequel les \omterm{def:g11:seismic-structure:body-waves}{ondes P} se propagent
plus lentement que dans le manteau.
\end{solution}

\begin{solution}{exo:g11:seismic-structure:12}
(1) D'après la proposition, $t_i = 2h\cos i_c/v_1$, et $\cos i_c = \sqrt{1 -
v_1^2/v_2^2} = \sqrt{v_2^2 - v_1^2}/v_2$. Donc
\[
  t_i = \frac{2h\sqrt{v_2^2 - v_1^2}}{v_1 v_2},
  \qquad\text{d'où}\qquad
  h = \frac{t_i\,v_1 v_2}{2\sqrt{v_2^2 - v_1^2}}.
\]
(2) $h = 6.23 \times 6.2 \times 8.1/\bigl(2\sqrt{65.61 -
38.44}\bigr) \approx 312.9/10.43 \approx \qty{30.0}{km}$.
\end{solution}

\begin{solution}{exo:g11:seismic-structure:13}
(1) $(3480/6371)^3 \approx 0.163$~: le noyau représente environ
\qty{16}{\percent} du volume. (2) La masse volumique moyenne est la moyenne
pondérée par les volumes~: $5.52 = 0.163
\rho_c + 0.837 \times 4.45$, donc $\rho_c \approx (5.52 - 3.72)/0.163 \approx
\qty{11.0}{g/cm^3}$. (3) C'est plus que \qty{9.9}{g/cm^3}, comme il se
doit, puisque la masse volumique du noyau augmente vers le centre (jusqu'à
\qty{13.1}{g/cm^3})~; c'est environ quatre fois la masse volumique de la
croûte, \qty{2.6}{g/cm^3}. Le noyau est fait d'une matière bien plus dense
que les roches de la surface.
\end{solution}

\begin{solution}{exo:g11:seismic-structure:14}
(1) $105 \times 111.2 \approx \qty{11700}{km}$ et $143 \times 111.2 \approx
\qty{15900}{km}$. (2) $(\cos\ang{105} - \cos\ang{143})/2 \approx (-0.259 +
0.799)/2 \approx 0.27$~: environ \qty{27}{\percent} de la surface, un anneau
autour du point opposé au séisme. (3) La \omterm{def:g11:seismic-structure:shadow-zone}{zone d'ombre} est un intervalle de
distances à l'épicentre~: elle se déplace avec le séisme. Chaque séisme a
son propre anneau.
\end{solution}

\begin{solution}{exo:g11:seismic-structure:15}
(1) $h = 9.5 \times 6.3 \times 8.1/\bigl(2\sqrt{65.61 - 39.69}\bigr) \approx
484.8/10.18 \approx \qty{48}{km}$. (2) $\sin i_c = 6.3/8.1 \approx 0.778$,
$\cos i_c \approx 0.629$, et $t_i = 2 \times 30 \times 0.629/6.3 \approx
\qty{6.0}{s}$. (3) Sous la chaîne, la croûte est épaisse d'environ
\qty{48}{km}, plus que la moyenne continentale d'environ \qty{41}{km} et
près de la moyenne des chaînes de montagnes~: les hautes chaînes reposent
sur une croûte épaisse.
\end{solution}

\begin{solution}{pb:g11:seismic-structure:1}
\textbf{1.} Les points Pg sont alignés sur une droite~; les points Pn sur une
autre droite, moins pentue.

\textbf{2.} Le séisme se produit à très faible profondeur sous l'épicentre~:
une station placée à distance nulle enregistrerait l'onde aussitôt.

\textbf{3.} $(32.3 - 3.3)/(200 - 20) = 29.0/180 \approx \num{0.161}$~\unit{s}
par kilomètre.

\textbf{4.} $v_1 = 180/29.0 \approx \qty{6.2}{km/s}$.

\textbf{5.} $3.3 + 280 \times 0.161 \approx \qty{48.4}{s}$.

\textbf{6.} $(43.2 - 18.6)/(300 - 100) = 24.6/200 = \num{0.123}$~\unit{s}
par kilomètre, donc $v_2 = 200/24.6 \approx \qty{8.13}{km/s}$.

\textbf{7.} Elle est proche de \qty{8.0}{km/s}~: sous la croûte se trouve la
roche du manteau le plus superficiel.

\textbf{8.} $\sin i_c = 6.21/8.13 \approx 0.763$, donc $i_c \approx
\ang{49.8}$.

\textbf{9.} L'\omterm{ex:g11:seismic-structure:sp}{onde conique} descend sous l'angle $i_c$ et remonte sous le
même angle. Le premier point où elle peut revenir à la surface se trouve à
$h\tan i_c$ au-delà du point où elle a atteint le \omterm{def:g11:seismic-structure:moho}{Moho}, lui-même à $h\tan i_c$
de la source~: aucune station plus proche que $2h\tan i_c$ ne peut la
recevoir.

\textbf{10.} Un sismomètre enregistre toutes les ondes qui l'atteignent, pas
seulement la première~; l'onde Pn arrive après l'onde Pg, mais elle se voit
encore comme une onde distincte sur l'enregistrement.

\textbf{11.} À \qty{160}{km}, Pg (\qty{25.8}{s}) arrive avant Pn
(\qty{26.0}{s})~; à \qty{180}{km}, Pn (\qty{28.5}{s}) arrive avant Pg
(\qty{29.0}{s}). Le changement se produit entre ces deux stations.

\textbf{12.} Pg~: $t = 3.3 + 0.1611\,(x - 20) \approx 0.08 + 0.1611\,x$. Pn~:
$t = 18.6 + 0.123\,(x - 100) = 6.30 + 0.123\,x$ ($x$ en kilomètres, $t$ en
secondes).

\textbf{13.} $0.08 + 0.1611\,x = 6.30 + 0.123\,x$ donne $x \approx
6.22/0.0381 \approx \qty{163}{km}$.

\textbf{14.} Pn, au bout de $6.30 + 0.123 \times 250 \approx \qty{37.1}{s}$.

\textbf{15.} $t_i \approx \qty{6.3}{s}$.

\textbf{16.} $\cos i_c \approx 0.646$, donc $h = 6.21 \times 6.3/(2 \times
0.646) \approx \qty{30.3}{km}$.

\textbf{17.} $x_c = 2 \times 30.3 \times \sqrt{14.34/1.92} \approx 60.6 \times
2.73 \approx \qty{165}{km}$, à \qty{2}{km} près des \qty{163}{km}
de la question~13.

\textbf{18.} $\tan i_c \approx 1.18$~: $2 \times 30.3 \times 1.18 \approx
\qty{72}{km}$. La station la plus proche qui enregistre Pn est à
\qty{80}{km}, et celle de \qty{60}{km} ne l'enregistre pas~: les deux
concordent.

\textbf{19.} Les temps sont lus à \qty{0.1}{s} près, et une erreur de
\qty{0.1}{s} sur le seul $t_i$ modifie $h$ d'environ \qty{0.5}{km}~; les
vitesses, tirées de pentes, sont elles aussi incertaines, et le séisme ne
s'est pas produit exactement en surface. La croûte n'est pas non plus
parfaitement uniforme. Une précision d'un ou deux kilomètres est le mieux
que permettent les données.

\textbf{20.} \textbf{Le \omterm{def:g11:seismic-structure:moho}{Moho} se trouve à environ \qty{30}{km} sous cette
région}~: la croûte y est plutôt mince pour un continent, où elle a en
moyenne une épaisseur d'environ \qty{41}{km}~; \qty{30}{km} est une valeur
typique des régions où la croûte a été étirée.
\end{solution}
